\clearpage
\chapter*{KATA PENGANTAR}
\phantomsection
\addcontentsline{toc}{chapter}{KATA PENGANTAR}
\setstretch{1.5}

Puji syukur penulis panjatkan kepada Tuhan Yang Maha Esa atas rahmat dan
karunia-Nya sehingga Laporan Tugas Besar Mata Kuliah \MataKuliah\ yang
berjudul ``\JudulTA'' dapat diselesaikan. Laporan ini disusun sebagai salah
satu syarat penyelesaian Tugas Besar pada Program Studi \ProgramStudi,
\Fakultas, \Universitas.

Sistem yang dikembangkan merupakan sistem monitoring suhu berbasis IoT
cerdas yang menggabungkan perangkat keras simulasi ESP32-S3, sensor DS18B20,
platform \textit{cloud} ThingSpeak, dan kecerdasan buatan berbasis
\textit{Random Forest Regressor}. Sistem dirancang secara \textit{end-to-end}
mulai dari akuisisi data sensor, komunikasi IoT, hingga prediksi suhu dan
visualisasi dashboard.

Penulis menyampaikan terima kasih kepada:
\begin{enumerate}[leftmargin=1.25cm]
    \item Tuhan Yang Maha Esa atas kelancaran dan kemudahan dalam penyelesaian
    tugas besar ini;
    \item orang tua dan keluarga yang senantiasa memberikan doa dan dukungan;
    \item \DosenPengampu, selaku Dosen Pengampu Mata Kuliah \MataKuliah; dan
    \item seluruh rekan mahasiswa Program Studi \ProgramStudi\ yang telah
    memberikan masukan dan semangat selama pengerjaan tugas besar ini.
\end{enumerate}

Penulis menyadari bahwa laporan ini masih dapat disempurnakan. Kritik dan
saran yang membangun sangat diharapkan. Penulis berharap sistem yang
dikembangkan dapat memberikan manfaat sebagai referensi implementasi IoT
cerdas berbasis simulasi dengan integrasi kecerdasan buatan.

\begin{flushright}
\Kota, \TanggalPengumpulan\\[1.5cm]
\NamaMahasiswa
\end{flushright}
